% ===================================================================
%  BAGAN No. 14 — Jalan. --- Para pedagang.
%
%  Traduction, faite en 2026, en indonesien d'aujourd'hui, sur l'ido
%  de texto/io. Regles de la colonne : voir 10-bagan-01.tex. Une seule
%  numerotation.
%
%  toko (15, 77), LA BOUTIQUE, EST HOKKIEN — ET C'EST LE MOT CENTRAL
%  D'UN TABLEAU QUI NE PARLE QUE DE BOUTIQUES. Il vient de 土庫, le
%  magasin, et il a supplante le kedai malais partout ou il s'agit
%  d'un commerce a devanture. La raison est historique et non
%  linguistique : ce sont les marchands chinois qui ont tenu les
%  boutiques de l'archipel pendant des siecles, et la langue a pris
%  leur mot pour la chose qu'ils tenaient. Le tableau 5 avait revele
%  cette couche par la peinture et le grenier ; ici elle occupe le
%  sujet meme du tableau.
%
%  ET LE MOULE « tukang » TIENT ENCORE. Le tableau 5 avait releve
%  seize metiers batis sur ce seul mot ; la rue en ajoute six : tukang
%  sapu le balayeur (11), tukang kaleng le ferblantier (14), tukang
%  kereta le carrossier (17), tukang pangkas le coiffeur (23), tukang
%  celup le teinturier (53), tukang topi le chapelier (39). Vingt-deux
%  metiers sur un seul moule dans une meme colonne, et pas un emprunt
%  pour nommer l'homme.
%
%  MAIS LES GRADES SONT NEERLANDAIS OU FRANCAIS, ET ILS CONFIRMENT
%  L'OURDOU. kolonel, mayor, kapten, letnan, batalyon, kompi,
%  resimen, jenderal, ajudan : c'est la meme liste qu'au tableau 14
%  ourdou, avec l'autre coloniser. Le releve a maintenant vu ce
%  partage sur trois administrations — britannique, neerlandaise et
%  celle de la marine — et il ne varie jamais : ce qui porte un
%  insigne s'emprunte a qui l'a cousu.
%
%  ET pesuruh REVIENT (79) : le meme mot qu'au tableau 13 pour le
%  chasseur d'hotel, employe ici pour le livreur en livree. Une langue
%  qui n'a pas de nom de fonction pour l'enfant qu'on envoie fait
%  servir le meme mot dans les deux tableaux — et le releve peut
%  maintenant dire que ce n'est pas un pis-aller de traduction mais un
%  fait de la langue.
% ===================================================================

%%K t14-apar-1 apar
\VUsaut{4.0mm}
\VUtitre{15.0pt}{60}{BAGAN No. 14}
\VUsaut{3.0mm}

%%K t14-tit-1 sub
\VUpk{11.4pt}{40}{Jalan. --- Para pedagang.}
\VUsaut{3.0mm}

%%K t14-01-1 p
1. --- Kawan saya Ioannes, yang tinggal di desa, datang menghabiskan
tiga hari di rumah keluarga saya. \VUgras{Sampai saat itu} ia belum
pernah berkesempatan melihat kota besar ; karena itu kami memakai
seluruh hari kami untuk berjalan-jalan, dan saya menjelaskan kepadanya
apa yang belum dikenalnya.

%%K t14-02-1 p
2. --- Mula-mula kami pergi mengunjungi \VUgras{observatorium}
\textsuperscript{(1)}, yang sangat menarik hatinya. Ketika kembali
lewat \VUgras{jalan} \textsuperscript{(3)} tempat \VUgras{pemandian
Turki} \textsuperscript{(2)} berada, kami berjumpa dengan sebuah
\VUgras{iring-iringan jenazah} \textsuperscript{(4)}. Di depan
\VUgras{iringan} itu berjalan \VUgras{para pendeta,} mendahului
\VUgras{kereta jenazah} tempat \VUgras{peti mati} diletakkan. Banyak
teman mengiringi keluarga yang \VUgras{berkabung.}

%%K t14-02-2 p
Sesudah rombongan itu lewat, kami menyeberangi jalan di depan
\VUgras{kedai bir} \textsuperscript{(5)} yang besar, tempat banyak
\VUgras{tamu} minum \VUgras{bir} di terasnya, terlindung oleh
\VUgras{tenda kaca} \textsuperscript{(6)}.

%%K t14-03-1 p
3. --- Ioannes sangat heran melihat \VUgras{perisai lambang} yang
dipasang di antara toko \VUgras{penjual minuman keras}
\textsuperscript{(7)} dan toko \VUgras{penjual alat musik}
\textsuperscript{(10)}. Ia menanyakan artinya kepada saya ; saya
menjawab bahwa perisai itu menandai \VUgras{konsulat}
\textsuperscript{(8)} Inggris, dan saya menunjukkan kepadanya
\VUgras{konsul} \textsuperscript{(9)} yang berdiri di balkon.

%%K t14-tit-2 sub
\VUpk{10.2pt}{40}{Melihat-lihat toko.}
\VUsaut{3.0mm}

%%K t14-04-1 p
4. --- Tentu saja kami berhenti di depan pajangan \VUgras{alat-alat
musik,} \VUgras{harmonium,} \VUgras{orgel} dan \VUgras{piano} ; kami
memandang sampul bergambar \VUgras{partitur} dan \VUgras{lagu-lagu}
untuk piano, dan mengagumi \VUgras{kecapi} dari kayu bersepuh emas yang
dipakai sebagai lambang.

%%K t14-05-1 p
5. --- Awan debu yang ditimbulkan \VUgras{para tukang sapu}
\textsuperscript{(11)} dan \VUgras{kereta penyapu}
\textsuperscript{(12)} memaksa kami cepat-cepat lewat di depan
\VUgras{perabot} indah milik \VUgras{penjual perabot}
\textsuperscript{(13)} dan di depan pajangan \VUgras{tungku} dan
\VUgras{lampu} milik \VUgras{tukang kaleng} \textsuperscript{(14)}.

%%K t14-06-1 p
6. --- Lagi pula, di tempat ini kami terpaksa meninggalkan trotoar,
sebab trotoar itu penuh dengan bungkusan yang diturunkan dari sebuah
\VUgras{kereta angkut} \textsuperscript{(16)} untuk diletakkan di dalam
sebuah \VUgras{toko} \textsuperscript{(15)} yang baru saja disewa.

%%K t14-06-2 p
Hal itu menghalangi kami melihat kereta-kereta indah milik
\VUgras{tukang kereta} \textsuperscript{(17)}, tetapi tidak menghalangi
kami berhenti memandang \VUgras{tukang kemas} \textsuperscript{(19)}
yang sedang membuat peti untuk tetangganya, \VUgras{penjual sepeda dan
mobil} \textsuperscript{(18)}. Kemudian, karena hari sudah agak
sore, kami pulang ke rumah.

%%K t14-tit-3 sub
\VUpk{10.2pt}{40}{Berjalan-jalan di kota.}
\VUsaut{3.0mm}

%%K t14-07-1 p
7. --- Keesokan harinya ibu saya mengusulkan agar Ioannes dipotret,
lalu menugasi saya menemaninya ke \VUgras{juru potret}
\textsuperscript{(20)}. Sesudah makan siang kami masuk ke
\VUgras{studio} seniman itu, tempat kami terpaksa menunggu cukup lama.
Untuk mengisi waktu kami memandang \VUgras{potret-potret}
\textsuperscript{(22)} yang tergantung di dinding.

%%K t14-07-2 p
Ketika tiba giliran kami, pekerjaan itu tidak berlangsung lama, dan
begitu kawan saya selesai berpose di depan \VUgras{kamera}
\textsuperscript{(21)}, kami turun.

%%K t14-08-1 p
8. --- Di lantai satu kami melihat lewat pintu \VUgras{tukang pangkas}
\textsuperscript{(23)}, yang terbuka, pembantunya sedang memotong
rambut seorang pelanggan.

%%K t14-08-2 p
Setiba di jalan, kami lewat di depan \VUgras{toko tembakau}
\textsuperscript{(24)}. Ioannes begitu terhibur oleh \VUgras{lentera}
\textsuperscript{(25)} merah dan \VUgras{pipa besar} yang dipakai
sebagai \VUgras{lambang} \textsuperscript{(26)} sehingga ia menabrak dua
tuan tua yang sedang berbincang sambil \VUgras{mengisap tembakau
hidung} \textsuperscript{(27)}.

%%K t14-tit-4 sub
\VUpk{10.2pt}{40}{Toko kue.}
\VUsaut{3.0mm}

%%K t14-09-1 p
9. --- \VUgras{Uang kertas} dan \VUgras{uang logam} dari segala negeri
yang dipamerkan di \VUgras{etalase} \VUgras{penukar uang}
\textsuperscript{(29)} menarik perhatian kami beberapa saat, tetapi
karena sudah waktunya makan kecil, kami masuk ke tempat \VUgras{tukang
kue} \textsuperscript{(31)} tanpa berhenti lebih lama.

%%K t14-10-1 p
10. --- Melihat segala \VUgras{penganan} yang ada di toko itu —
\VUgras{kue, kue madu, biskuit} dan segala macam \VUgras{gula-gula} —
kami sulit sekali memilih. Akhirnya Ioannes memilih \VUgras{kue krim,}
dan saya hanya mengambil sepotong kecil \VUgras{tar ceri,} sebab
makanan manis \VUgras{membuat gigi saya sakit,} dan saya sama sekali
tidak ingin terlalu sering pergi ke \VUgras{dokter gigi}
\textsuperscript{(30)}.

%%K t14-tit-5 sub
\VUpk{10.2pt}{40}{Mengetik dengan mesin.}
\VUsaut{3.0mm}

%%K t14-11-1 p
11. --- Untuk menunjukkan kepada kawan saya \VUgras{mesin tik} yang
pernah didengarnya tetapi belum dikenalnya, kami masuk ke
\VUgras{kantor} \textsuperscript{(32)} \VUgras{sepupu laki-laki saya}
\textsuperscript{(34)}. Kami mendapatinya sedang mendiktekan
surat-suratnya kepada \VUgras{sekretarisnya} \textsuperscript{(35)},
yang menjadi \VUgras{juru steno.} Baru saja surat itu selesai, seorang
\VUgras{juru ketik} \textsuperscript{(33)} menyalinnya dengan mesinnya.
Karena tidak ingin mengganggu pekerjaan kerabat saya, kami hanya
tinggal beberapa saat di sana.

%%K t14-12-1 p
12. --- Di bawah kantor sepupu saya tinggal seorang \VUgras{tukang jam
dan pandai emas} \textsuperscript{(37)} yang besar, yang juga menjadi
\VUgras{pandai emas.} Tokonya yang megah berisi banyak \VUgras{jam,
jam bandul, arloji saku,} \VUgras{rantai} dan \VUgras{perhiasan} yang
berharga — misalnya \VUgras{kalung mutiara,} \VUgras{gelang} emas,
\VUgras{anting} dan \VUgras{cincin} yang bertatah \VUgras{batu permata}
dan \VUgras{berlian.} Salah satu etalase khusus disediakan untuk
\VUgras{barang emas} dan berisi koleksi \VUgras{alat makan} perak yang
besar.

%%K t14-13-1 p
13. --- Ketika berbelok di sudut jalan, saya melihat bahwa
\VUgras{papan nama} \textsuperscript{(36)} yang dipasang di dekat
\VUgras{nomor} rumah telah diganti. Saya hampir mengatakannya
keras-keras ketika terdengar bunyi riang \VUgras{korps} musik tentara.
Kami berlari menyongsong resimen itu, tanpa memikirkan bahwa
\VUgras{ia} akan lewat di depan toko \VUgras{tukang topi}
\textsuperscript{(39)} dan \VUgras{tukang sarung tangan}
\textsuperscript{(40)}, dan bahwa kami akan sama baiknya menyaksikan
pawainya bila tetap berdiri di depan etalase \VUgras{tukang kacamata}
\textsuperscript{(38)}, yang penuh dengan \VUgras{kacamata jepit,
kacamata bertangkai} dan \VUgras{teropong kecil.}

%%K t14-tit-6 sub
\VUpk{10.2pt}{40}{Pawai.}
\VUsaut{3.0mm}

%%K t14-14-1 p
14. --- Di depan \VUgras{resimen} \textsuperscript{(41)} berjalan
\VUgras{kepala tambur} \textsuperscript{(42)}, diikuti \VUgras{para
penabuh tambur} \textsuperscript{(43)}, \VUgras{para peniup sangkakala}
\textsuperscript{(44)} dan seluruh \VUgras{korps musik.}
\VUgras{Jenderal} \textsuperscript{(45)}, yang berada di alun-alun
bersama ajudannya, berhenti di dekat \VUgras{pancuran air}
\textsuperscript{(49)} \VUgras{air mancur} \textsuperscript{(48)} untuk
menyaksikan \VUgras{pawai.}

%%K t14-14-2 p
\VUgras{Para perwira staf} \textsuperscript{(46)}, \VUgras{kolonel} dan
\VUgras{letnan kolonel} memberi hormat kepadanya dengan pedang.
\VUgras{Para mayor} berada di depan \VUgras{batalyon}
\textsuperscript{(47)} mereka dan setiap \VUgras{kapten} memimpin
\VUgras{kompinya,} sementara \VUgras{para letnan,} \VUgras{letnan muda}
dan \VUgras{bintara} menempati tempat mereka yang biasa di samping
anak buahnya.

%%K t14-15-1 p
15. --- Banyak orang yang lewat berkumpul di \VUgras{persimpangan}
\textsuperscript{(50)} ; para pedagang — \VUgras{tukang payung}
\textsuperscript{(51)}, \VUgras{tukang celup} \textsuperscript{(53)},
\VUgras{tukang senjata} \textsuperscript{(54)} — keluar untuk melihat
\VUgras{para prajurit.} Semua kendaraan — \VUgras{kereta sewa}
\textsuperscript{(52)}, \VUgras{gerobak panjang} \textsuperscript{(57)}
dan juga \VUgras{tong penyiram} \textsuperscript{(56)} — berhenti di
sepanjang trotoar.

%%K t14-tit-7 sub
\VUpk{10.2pt}{40}{Pemandangan di jalan.}
\VUsaut{3.0mm}

%%K t14-15-2 p
Seorang \VUgras{pedagang keliling} \textsuperscript{(55)}, yang
menjinjing \VUgras{kotak contohnya} dengan tangan, dengan cepat
menyeberangi jalan, dan orang-orang yang duduk di \VUgras{tingkat atas}
\textsuperscript{(59)} \VUgras{omnibus} \textsuperscript{(58)} berbalik
untuk menikmati tontonan itu. Seorang \VUgras{penyanyi jalanan}
\textsuperscript{(60)} yang malang dengan \VUgras{organ putarnya}
\textsuperscript{(61)} terpaksa menghentikan \VUgras{lagunya,} sebab
gemuruh tambur menenggelamkan suaranya.

%%K t14-16-1 p
16. --- Ketika resimen itu sampai di depan \VUgras{toko kain}
\textsuperscript{(64)}, \VUgras{para penjahit perempuan}
\textsuperscript{(63)} dan \VUgras{para penjahit} \textsuperscript{(62)}
meninggalkan \VUgras{mesin jahit} dan pekerjaan mereka untuk memandang
lewat jendela. \VUgras{Pedagang} \textsuperscript{(65)}, yang baru saja
memasang \VUgras{setelan-setelan} \textsuperscript{(66)} baru di
\VUgras{etalasenya,} terpaksa menyuruh memasukkan ke dalam
\VUgras{kain-kain wol} \textsuperscript{(67)} dan \VUgras{kain katun}
yang dipasang di trotoar, di dekat \VUgras{lubang got}
\textsuperscript{(68)}, karena takut kerumunan akan merobohkannya ;
bahkan \VUgras{pedagang asongan} \textsuperscript{(69)} tua, yang
sedang ditawar kaus kakinya oleh seorang perempuan penjaga pintu,
terpaksa masuk ke sebuah lorong untuk menyelesaikan jualannya.

%%K t14-16-2 p
Kami mengiringi resimen itu sampai ke barak, \VUgras{dan,} dengan
sangat puas atas hari kami, kami pulang pada waktu makan malam.

%%K t14-tit-8 sub
\VUpk{10.2pt}{40}{Bermacam-macam usaha dan pekerjaan.}
\VUsaut{3.0mm}

%%K t14-17-1 p
17. --- Keesokan harinya ayah saya mengusulkan agar kami mengunjungi
percetakan koran setempat. Kami menerimanya dengan gembira.

%%K t14-17-2 p
\VUgras{Pencetak} itu juga \VUgras{penjual buku} \textsuperscript{(70)}
\VUgras{(yaitu yang menjual buku),} \VUgras{penerbit,} \VUgras{penjual
kertas} dan \VUgras{penjilid.} Ia memasang di lantai satu
\VUgras{ruang redaksi} \textsuperscript{(71)} koran itu, dan juga
\VUgras{ruang ukir,} tempat \VUgras{para tukang litografi}
\textsuperscript{(72)} dan \VUgras{para pengukir tembaga}
\textsuperscript{(73)} bekerja ; akhirnya \VUgras{ruang susun huruf,}
tempat \VUgras{para pekerja cetak} \textsuperscript{(74)} berada.
\VUgras{Ruang cetak} \textsuperscript{(75)} yang sesungguhnya,
\VUgras{mesin-mesin cetak} dan bermacam-macam mesin lain berada di
lantai dasar.

%%K t14-tit-9 sub
\VUpk{10.2pt}{40}{Ioannes bertanya banyak sekali.}
\VUsaut{3.0mm}

%%K t14-18-1 p
18. --- Keluar dari percetakan, Ioannes berhenti dengan sangat heran
di depan sebuah kotak kaca kecil yang dipasang pada sebuah tiang, di
seberang \VUgras{toko kelontong} \textsuperscript{(77)}, dan menanyakan
untuk apa kotak itu. Ayah saya menjelaskan kepadanya bahwa itu
\VUgras{kotak alarm kebakaran} \textsuperscript{(76)}. Lagi pula segala
sesuatu di kota itu mengherankan kawan saya, dari \VUgras{air mancur
batu} \textsuperscript{(78)} yang sederhana dan \VUgras{roda tiga
pengangkut} \textsuperscript{(80)} yang ringan, yang dikemudikan
seorang \VUgras{pesuruh} \textsuperscript{(79)} berseragam, sampai
\VUgras{kereta pos} \textsuperscript{(81)}.

%%K t14-19-1 p
19. --- Sementara ayah saya meninggalkan kami sejenak untuk berbicara
dengan \VUgras{penagih pajak} \textsuperscript{(83)}, yang berhenti
untuk berbincang dengan seorang \VUgras{pengacara}
\textsuperscript{(82)} dan seorang \VUgras{notaris}
\textsuperscript{(84)}, kami menghibur diri dengan mengikuti lalu
lalang di jalan dengan mata. Orang yang paling bermacam-macam lewat di
depan kami : seorang \VUgras{pesuruh tukang kue}
\textsuperscript{(85)}, yang membawa \VUgras{pastel} yang megah di
dalam keranjang, datang di belakang seorang \VUgras{pedagang pakaian}
\textsuperscript{(86)} ; seorang \VUgras{penyala lentera}
\textsuperscript{(87)} sedang berbincang dengan seorang \VUgras{petugas
gas} \textsuperscript{(88)} ; seorang \VUgras{biarawati}
\textsuperscript{(89)} yang menggandeng seorang anak yatim sedang
kembali ke \VUgras{biaranya.}

%%K t14-tit-10 sub
\VUpk{10.2pt}{40}{Dalam perjalanan pulang.}
\VUsaut{3.0mm}

%%K t14-20-1 p
20. --- Pada saat ayah saya menyusul kami, Ioannes tertawa
terbahak-bahak melihat wajah kotor dan pakaian koyak yang aneh dua
\VUgras{tukang sapu cerobong} \textsuperscript{(91)} yang penuh jelaga.

%%K t14-21-1 p
21. --- Saya hendak membeli sebatang tanaman dari \VUgras{penjual
bunga} \textsuperscript{(92)} untuk ibu saya, tetapi ayah saya
mengingatkan bahwa tanaman itu akan terlalu berat dibawa dan lebih
baik membeli \VUgras{jeruk.} Kami mendekati \VUgras{penjual perempuan}
\textsuperscript{(94)}, dan ketika \VUgras{pengasuh}
\textsuperscript{(93)}, yang sedang memilih untuk anak asuhnya,
selesai berbelanja, kami pun dilayani.

%%K t14-22-1 p
22. --- Dalam perjalanan pulang kami berjumpa dengan beberapa kenalan
: seorang teman lama keluarga saya, yang bertumpu pada lengan
\VUgras{nyonya pendampingnya} \textsuperscript{(95)} ; \VUgras{insinyur}
\textsuperscript{(96)} jembatan dan jalan ; dan salah seorang sepupu
perempuan saya bersama \VUgras{pengasuh anaknya}
\textsuperscript{(98)}.

%%K t14-22-2 p
Kemudian kami berjumpa dengan \VUgras{tukang topi perempuan}
\textsuperscript{(97)} ibu saya dan dua \VUgras{biarawan}
\textsuperscript{(99)} asing bagi kota itu, yang menghampiri kami untuk
menanyakan jalan ke stasiun kepada ayah saya.

\VUsaut{6.0mm}
\VUfilet{22mm}
